\documentclass[review,times]{elsarticle}
\usepackage{amsmath,amssymb,amsthm,mathtools}
\usepackage{graphicx}
\usepackage{booktabs}
\usepackage{lineno}
\linenumbers

\newtheorem{theorem}{Theorem}
\newtheorem{lemma}{Lemma}
\newtheorem{corollary}{Corollary}
\theoremstyle{definition}
\newtheorem{definition}{Definition}
\newtheorem{assumption}{Assumption}
\theoremstyle{remark}
\newtheorem{remark}{Remark}
\DeclareMathOperator{\dtv}{d_{\mathrm{TV}}}
\newcommand{\law}[1]{\mathcal{V}_{#1}}
\newcommand{\lawt}[1]{\tilde{\mathcal{V}}_{#1}}
\journal{Reliability Engineering \& System Safety}

\begin{document}
\begin{frontmatter}
\title{Similarity-Calibrated Conformal prediction: data-free coverage guarantees for
remaining-useful-life intervals under operating-regime transfer}
\author{Bien Busico}
\address{Map\'ua Malayan College Mindanao, Davao City, Philippines}

\begin{abstract}
In condition-based maintenance and safety decisions, a remaining-useful-life (RUL) interval is only
as good as its stated confidence. Conformal prediction gives a
distribution-free coverage guarantee, but only under exchangeability, which fails when a model
calibrated on assets in one operating regime is deployed in another: the interval then miscovers
with no warning. The standard fixes do not suit reliability practice: reweighting needs target
failure data a new asset lacks, and dimensionless transfer learning carries no coverage guarantee.
We present Similarity-Calibrated Conformal (SCC) prediction, which calibrates the
conformal quantile in a physics-derived dimensionless coordinate (via the Buckingham $\Pi$ theorem)
so coverage transfers across operating regimes with no target failure data. On a controlled
catalyst-deactivation testbed, regime-blind conformal miscovers (coverage falls to 0.572) while SCC
recovers it to within 0.011 of target, degrades monotonically under a graded departure from
similitude and stays actionable to a quantified threshold. The calibration replicates on 509
externally authored turbofan engines, where the gap falls from 0.590 to 0.060.
SCC adds two features a practitioner can act on: a runtime diagnostic returning holds, violated or
indeterminate, so the method declines to certify rather than miscover; and an a-priori coverage
bound, at most $2L\lVert\Delta\boldsymbol\psi\rVert$ in a similitude-departure quantity computed
from physical parameters. The bound is valid but conservative, and we claim it only where the
governing degradation law identifies $\boldsymbol\psi$, as ISO 281 does for rolling-contact fatigue
and no comparable standard does for a turbofan gas path. Validation on adequately powered
operational field data is the next step.
\end{abstract}

\begin{keyword}
conformal prediction \sep uncertainty quantification \sep remaining useful life \sep
dimensional analysis \sep transfer learning \sep prognostics
\end{keyword}
\end{frontmatter}

% ---------- begin sections/01_introduction ----------
\section{Introduction}\label{sec:intro}
Condition-based maintenance and safety-critical operation both rest on prognostics: an estimate of
an asset's remaining useful life (RUL) that a planner can act on. What a maintenance or replacement
decision actually needs is not a point estimate but a trustworthy \emph{lower bound} on remaining
life, with a stated confidence that the true life will not fall below it. An interval that is too
optimistic invites in-service failure, with the attendant safety and cost consequences; one that is
too conservative forces premature replacement and lost availability. The usefulness of a prognostic
interval is therefore set by whether its stated confidence is honest.

Conformal prediction has become an attractive way to attach such a guarantee to a data-driven RUL
model, because it certifies a finite-sample, distribution-free coverage level without assuming a
noise model~\cite{vovk2005algorithmic,papadopoulos2002inductive,lei2018distribution,angelopoulos2023gentle,javanmardi2023conformal}.
That guarantee, however, holds only when the calibration and deployment data are exchangeable, and
this is precisely the assumption that fails in the situation reliability engineering meets most
often: a model calibrated on assets run under one operating regime (load, speed, temperature) is
deployed on assets run under another. The degradation distribution shifts, exchangeability is lost,
and the interval miscovers with no outward sign that its confidence is no longer valid.

The two established ways to repair coverage under such a shift do not fit the reliability setting.
Reweighting methods, namely covariate-shift conformal prediction~\cite{tibshirani2019conformal} and
its non-exchangeable generalisation~\cite{barber2023conformal}, restore coverage, but only by
estimating the likelihood ratio between the target and source operating distributions, which
requires representative failure data from the target regime. A newly commissioned asset, or a regime
with no run-to-failure history, does not have that data. Dimensional-analysis transfer, in which a
model is expressed in dimensionless variables so that it generalises across scales and
systems~\cite{oppenheimer2023multiscale,therrien2024buckingham,girard2024dimensionless}, needs no
target data, but it improves only the \emph{point} prediction and carries no coverage guarantee.
One family has the guarantee but needs data the operator lacks; the other needs no data but offers
no guarantee.

This paper closes that gap with a deployable method, Similarity-Calibrated Conformal (SCC)
prediction, summarised in Fig.~\ref{fig:overview}. SCC computes the conformal calibration in a
dimensionless coordinate chosen, from the governing physics, so that the score distribution is
invariant across operating regimes under dynamic similitude. Coverage then transfers with no target
failure data, and the residual shortfall when similitude is imperfect is bounded by a physical
quantity computable before deployment. A companion runtime check tells the operator when that
premise is supported, so the method declines to certify rather than miscover silently.

The contributions are practical capabilities for reliable prognostics under operating-regime
transfer:
\begin{enumerate}
\item a calibration step in dimensionless ($\Pi$) coordinates that restores conformal coverage
across operating regimes without target-regime failure data (Section~\ref{sec:method});
\item an a-priori coverage bound computable from physical parameters alone: the cross-regime
coverage shortfall is at most $2L\lVert\Delta\boldsymbol\psi\rVert$ in a similitude-departure
quantity $\boldsymbol\psi$, so the worst-case loss can be checked before an asset is deployed
(Section~\ref{sec:method}; derivation in~\ref{app:proof});
\item a runtime similitude diagnostic returning holds, violated, or indeterminate, which gives SCC a
self-aware operating envelope and a conservative fallback when its premise is unsupported
(Section~\ref{sec:method});
\item a non-circular validation on a controlled process-reliability testbed, stress-testing the
guarantee against a graded, physically meaningful departure from similitude, together with an
external replication on 509 run-to-failure turbofan engines we did not author, on which the
coverage recovery holds and the a-priori bound is reported to fail
(Sections~\ref{sec:design}--\ref{sec:results}).
\end{enumerate}

We are deliberate about scope. The guarantee is conditional on two physical hypotheses that are
falsifiable rather than assumed; the validation uses a controlled simulation in which similitude is
real but is never assumed by the estimator; the coverage recovery is replicated on an externally
authored turbofan fleet, where the a-priori bound is also checked and reported to fail; and the thin
FEMTO bearing benchmark (an accelerated laboratory test platform rather than in-service field data)
is reported by the diagnostic as outside the method's envelope rather than forced into a positive
result. Validation on adequately powered operational field data is the natural next step.
% ---------- end sections/01_introduction ----------
% ---------- begin sections/02_background ----------
\section{Background and related work}\label{sec:background}
\emph{Split conformal prediction} gives the guarantee SCC builds on. Given a nonconformity score
$V=s(X,Y)$ and a calibration set exchangeable with the test point, the prediction set formed from
the empirical $(1-\alpha)$ quantile of the calibration scores contains the truth with probability at
least $1-\alpha$~\cite{papadopoulos2002inductive,lei2018distribution}. The guarantee is marginal,
finite-sample, and free of modelling assumptions; its only structural requirement is
exchangeability, which is what operating-regime transfer destroys.

When exchangeability fails, Barber et al.~\cite{barber2023conformal} bound the resulting coverage
loss by the total-variation distances between the score vector and its versions with the test point
swapped against each calibration point; for independent data this reduces to pairwise
total-variation distances between the source and target score laws. This non-exchangeable bound is
the analytical backbone of the SCC certificate. Weighted (covariate-shift) conformal
prediction~\cite{tibshirani2019conformal} instead restores exact coverage using importance weights
$w(x)=\mathrm{d}P_T/\mathrm{d}P_S$; this is powerful but requires estimating the density ratio, and
hence representative target covariates. SCC differs in kind: it removes the shift physically, through
a dimensionless transformation, rather than reweighting statistically, and its certificate needs no
target data.

The physical transformation comes from \emph{dimensional similitude}. The Buckingham $\Pi$ theorem
reduces a dimensional relationship to dimensionless groups, and matching those groups enforces
dynamic similitude between systems. Recent work exploits this for transfer learning, training models
in dimensionless variables so they generalise across scales and operating
points~\cite{oppenheimer2023multiscale,therrien2024buckingham}, with exact policy transfer
guaranteed only under exact similitude~\cite{girard2024dimensionless}. These methods target point
accuracy and provide no calibrated uncertainty; SCC imports the same similitude principle into the
conformal calibration step to obtain a coverage guarantee.

The closest prior treatment of prognostic uncertainty is Javanmardi and
H\"ullermeier~\cite{javanmardi2023conformal}, who first conformalised RUL estimation on the C-MAPSS
benchmark using split conformal, conformalised quantile regression, and the non-exchangeable
variant. Where operating conditions vary, their remedy is to cluster them into a finite set of
\emph{known, discrete} modes and normalise within each; they note that this becomes inapplicable
when the modes are not known in advance or lack calibration data, and their non-exchangeable variant
weights calibration points by temporal recency rather than operating-regime similarity. SCC instead
handles \emph{continuous} operating-regime transfer through physical scaling, certifies coverage a
priori without target-regime data, and needs no discrete mode labels. Our concern is orthogonal to
point accuracy~\cite{lei2018machinery}: it is the validity of the interval when the operating regime
changes between calibration and deployment.
% ---------- end sections/02_background ----------
% ---------- begin sections/03_method ----------
\section{Method: Similarity-Calibrated Conformal prediction}\label{sec:method}
SCC is a drop-in modification of split conformal prediction that makes the calibration transferable
across operating regimes. It runs in two stages (Fig.~\ref{fig:overview}): offline, on a source asset
with run-to-failure data, it calibrates a conformal quantile in a dimensionless coordinate; online, on
a target asset in a different regime and with no failure history, it re-scales that quantile by the
target's physical scale to produce a certified RUL lower bound. The only quantities crossing from
source to target are the scale-free conformal quantile and the target's physical scale, so no target
failure data is needed.

\begin{figure}[t]\centering
\includegraphics[width=\linewidth]{fig_overview.pdf}
\caption{The SCC workflow. Offline, source scores are made dimensionless by a physics-derived scale
$\sigma(\theta_S)$ and a conformal quantile $\tilde q$ is taken. Online, the quantile is
re-dimensionalised by the target scale $\sigma(\theta_T)$ to give a certified RUL lower bound. The
coverage shortfall is bounded a priori from physical parameters, and a runtime diagnostic reports
whether the similitude premise is supported.}\label{fig:overview}
\end{figure}

\subsection{Setup and the SCC procedure}\label{sec:procedure}
Units are indexed by operating condition $c$ with known physical parameters $\theta_c$ (load, speed,
temperature). A unit yields features $X$ and true RUL $Y$ with joint law $\mathbb P_c$; a pre-trained
predictor $\hat y$ and the one-sided over-prediction score $V=\hat y(X)-Y$ (a lower RUL bound) have
score law $\law c$. The procedure is:
\begin{enumerate}
\item \emph{Dimensionless calibration.} On calibration data $\{(X_i,Y_i)\}_{i=1}^n$ drawn from the
source regime $S$, form scores $V_i=s(X_i,Y_i)$ and divide by the source scale to obtain dimensionless
scores $\tilde V_i=V_i/\sigma(\theta_S)$.
\item \emph{Conformal quantile.} Take $\tilde q=\tilde V_{(\lceil(1-\alpha)(n+1)\rceil)}$, the order
statistic giving the empirical $(1-\alpha)$ quantile of the dimensionless scores.
\item \emph{Re-dimensionalised deployment.} For a target unit under condition $T$, return the RUL set
$\mathcal C(X_{n+1})=\{y:s(X_{n+1},y)\le\sigma(\theta_T)\,\tilde q\}$, re-scaling the quantile by the
\emph{target} scale.
\end{enumerate}
Naive (regime-blind) conformal is the special case $\sigma\equiv1$; it applies the source quantile
directly to the target and miscovers whenever $\sigma(\theta_S)\neq\sigma(\theta_T)$. The single
change in SCC is calibrating and deploying in the $\sigma$-normalised coordinate.

\subsection{The dimensionless scale and what it absorbs}\label{sec:scale}
The scale $\sigma(\theta_c)$ comes from the governing physics through the Buckingham $\Pi$ theorem,
which yields dimensionless groups. These split into scale-setting groups, used to form a
characteristic score scale $\sigma(\theta_c)>0$ such that the dimensionless score
$\tilde V=V/\sigma(\theta_c)$ has a condition-invariant law under exact similitude, and a residual
\emph{similitude-departure} parameter $\boldsymbol\psi(\theta_c)\in\mathbb R^m$ collecting the groups
the scale does not absorb (unmodelled physics). When the dimensionless model is complete,
$\boldsymbol\psi$ is constant across conditions.

Two standard degradation laws show what the scale removes. For rolling-contact fatigue
(Lundberg--Palmgren / ISO~281~\cite{lundberg1947dynamic,iso281}, rating life $\propto(C/P)^p$ with
$p=3$), the scale $\sigma=T_L\propto(C/P)^p/n$ \emph{absorbs} the load--speed group $P/C$. For
first-order catalyst deactivation (Arrhenius $k=A\,e^{-E/RT}$~\cite{levenspiel1999chemical,perry2008handbook}),
the scale $\sigma=1/k_{\mathrm{nom}}$ \emph{contains} the primary rate group $e^{-E/RT}$. In both
cases a large difference in the dominant operating parameter is taken up by $\sigma$ and does not by
itself induce a mismatch in the score distribution; what remains is only the residual departure
$\boldsymbol\psi$.~\ref{app:bearing} carries the rolling-contact case through to a
numerical instance and identifies $\boldsymbol\psi$ for that law explicitly, which is where the abstraction
becomes an instruction a practitioner can act on.

\subsection{Coverage guarantee}\label{sec:guarantee}
Calibrating in the dimensionless coordinate restores coverage up to a term set by the physical
departure. With calibration drawn i.i.d.\ from the source $S$ and a target point
$(X_{n+1},Y_{n+1})\sim\mathbb P_T$ independent of calibration, the SCC interval of
Section~\ref{sec:procedure} satisfies
\begin{equation}\label{eq:cert}
\mathbb P\big(Y_{n+1}\in\mathcal C(X_{n+1})\big)\;\ge\;
1-\alpha-2L\lVert\boldsymbol\psi(\theta_S)-\boldsymbol\psi(\theta_T)\rVert,
\end{equation}
where $\lVert\Delta\boldsymbol\psi\rVert=\lVert\boldsymbol\psi(\theta_S)-\boldsymbol\psi(\theta_T)\rVert$
is the similitude departure between the regimes and $L$ is a Lipschitz constant estimated from source
data (Section~\ref{sec:estL}). Equation~\eqref{eq:cert} follows by composing the dimensionless
reduction with the non-exchangeable conformal bound of Barber et al.~\cite{barber2023conformal}; the
formal statement, hypotheses, and proof are collected in~\ref{app:proof}. Here we read off its
engineering content.

The departure $\lVert\Delta\boldsymbol\psi\rVert$ is a physical quantity, not a function of target
failure data, which gives two regimes of use. Under complete similitude ($\boldsymbol\psi$ constant)
the bound is zero and coverage is at least $1-\alpha$ for \emph{arbitrarily large} differences in the
scale-setting parameters, needing only the target scale $\sigma(\theta_T)$; this is the property a
newly commissioned asset can rely on. Under incomplete similitude, coverage is certified down to
$1-\alpha-2L\lVert\Delta\boldsymbol\psi\rVert$ once the departure is bounded, so the worst-case
shortfall is known before deployment. Neither case estimates a density ratio, which is what separates
SCC from weighted and covariate-shift conformal.

The dimension $m$ of $\boldsymbol\psi$ is set by the physics rather than by the method, and the case
exercised in this paper has $m=1$: the departure is the scalar weight on a single unmodelled
channel, so the norm in~\eqref{eq:cert} reduces to an absolute value and its choice is immaterial.
When $m>1$, as it becomes for rolling-contact fatigue once the contamination group is carried
alongside the lubrication regime (\ref{app:bearing}), the norm has to be fixed at the same time as
$\boldsymbol\psi$, since $L$ is estimated in that norm and rescales with it.

\subsection{What exchangeability requires, and block calibration}\label{sec:exch}
The calibration hypothesis behind~\eqref{eq:cert} is exchangeability across \emph{units}, not across
time steps within a unit. A monitored asset yields a score at every inspection point, and stacking
several points per unit inflates the nominal calibration size without adding a matching amount of
information: on the testbed of Section~\ref{sec:design}, scores taken at three monitoring fractions
of the same unit correlate at $\rho=0.908$ on average, a design effect $1+(F-1)\rho=2.82$, so 1200
stacked scores carry about 426 independent ones against 400 units. We therefore calibrate with one
score per unit throughout, and the reported calibration size is a count of units.

Where the guarantee wanted is coverage \emph{simultaneously at every monitoring point} of a unit
rather than at one sampled point, the principled generalisation is block (per-unit) conformal
calibration on the per-unit maximum score, which is exchangeable across units by construction. It
is the stronger maintenance statement and it costs conservatism: on the testbed it gives a coverage
gap of 0.007 at $\eta=0$ and 0.173 at $\eta=2$, against 0.011 and 0.097 for single-point unit-level
calibration. We report single-point unit-level calibration as the default and block calibration as
the option for multi-point monitoring.

\subsection{Estimating the departure sensitivity $L$}\label{sec:estL}
The constant $L$ in~\eqref{eq:cert} is obtained from source data alone. With at least two source
conditions, regress the empirical total-variation distance between dimensionless score laws,
$\dtv(\tilde{\mathcal V}_S,\tilde{\mathcal V}_{S'})$, on the physical departure
$\lVert\Delta\boldsymbol\psi\rVert$; the slope estimates $L$ and the intercept $a$ estimates the
finite-sample floor, which decays as $n^{-1/2}$. The design points are configurations, not conditions:
each ordered pair of source conditions contributes one point per distinct departure magnitude the
source data exhibits, so a fleet with three conditions supplies far more than three points. In the
study of Section~\ref{sec:res-cert} this gives 30 (pair, departure) configurations, of which 60\% are
used for the fit and the remaining 40\% are held out to validate the resulting bound. The operational bound
$2(a+L\lVert\Delta\boldsymbol\psi\rVert)$ then upper-bounds the target coverage gap using only
source-side information.

\subsection{Runtime similitude diagnostic}\label{sec:diagnostic}
The guarantee~\eqref{eq:cert} is conditional on the similitude premise, which bites only when
$\boldsymbol\psi$ varies across conditions. Because exact similitude makes the dimensionless life
condition-invariant, the premise is testable from source data, and SCC carries a runtime check that
acts as a go/no-go monitor. For each condition pair it forms
$g=\mathrm{median}(D\mid c)/\mathrm{median}(D\mid c')$, where $D$ is the dimensionless life (the
dynamic-capacity constant cancels), equivalently the empirical median-life ratio divided by the L10
prediction. A bootstrap confidence interval on $g$ yields a three-way verdict: $1$ inside a tight
interval gives \emph{holds}; $1$ outside gives \emph{violated}; an interval wider than a preset fold
gives \emph{indeterminate} (underpowered). On the latter two, SCC declines the a-priori certificate
and falls back to data-driven conformal, so an unsupported premise produces a conservative fallback
rather than a silent miscoverage.

The same check does double duty. Because it tests the invariance of the dimensionless life, it is
sensitive not only to a departure from similitude but also to misspecification of the scale itself:
an assumed activation energy that is wrong distorts the between-condition behaviour of $\sigma$ and
breaks the same invariance. Section~\ref{sec:res-misspec} reports it returning \emph{violated} at
every misspecification tested, including the smallest one at which a wrong scale makes SCC worse
than not scaling at all, so both failure modes are caught by one runtime test.
% ---------- end sections/03_method ----------
% ---------- begin sections/04_experimental_design ----------
\section{Case studies: a controlled testbed and an external turbofan fleet}\label{sec:design}
\subsection{Why a controlled testbed}\label{sec:why}
Demonstrating an a-priori certificate honestly requires data in which the governing physics actually
controls degradation, so that the similitude departure $\boldsymbol\psi$ is meaningful and the
certificate can be checked rather than asserted. We therefore use a stochastic catalyst-deactivation
simulation~\cite{levenspiel1999chemical,perry2008handbook}, a process-engineering reliability problem
in which a catalytic unit ages as its catalyst deactivates and reaches end of life when activity
falls below a threshold; the prognostic task is to predict the remaining life for catalyst-replacement
scheduling. Dynamic similitude is \emph{genuine} in the simulator but is never assumed by the
estimator, which sees only noisy activity signals and a nominal rate model. This separation permits a
non-circular test: we impose a graded, physically meaningful departure from similitude and check that
the certificate degrades exactly as~\eqref{eq:cert} predicts.

\subsection{Simulator}\label{sec:sim}
Deactivation follows first order, $\mathrm{d}a/\mathrm{d}t=-k_{\mathrm{eff}}(T)\,a$, with failure at
activity $a\le0.3$. The primary Arrhenius channel is $k_1=A_1e^{-E_1/RT}$ ($E_1=80$~kJ/mol); an
optional secondary channel $k_2=A_2e^{-E_2/RT}$ of weight $\eta$ provides the controlled departure
($E_2\neq E_1$). Unit-to-unit variability enters as a lognormal pre-exponential, and the activity
signal carries additive noise. At $\eta=0$ the dimensionless clock $\tau=t\,k_1(T)$ collapses every
temperature onto $a(\tau)=e^{-\tau}$: exact similitude, with the characteristic life $T_L=1/k_1(T)$
spanning $10.2\times$ across the three temperatures (600/650/700~K). The weight $\eta$ thus moves the
testbed continuously from exact similitude ($\eta=0$) into progressively stronger violations.

\subsection{Predictor, scores, and reliability metrics}\label{sec:metrics}
The predictor extrapolates remaining life from the observed activity using the nominal rate, and the
conformal layer covers the residual. Each metric is read in maintenance terms. \emph{Coverage}
(target $1-\alpha=0.90$) is the fraction of intervals that contain the true RUL, i.e.\ the realised
reliability of the maintenance decision; the \emph{coverage gap} $\max(0,(1-\alpha)-\text{coverage})$
is the shortfall against the promised reliability. \emph{Efficiency} is the relative interval
back-off $b=q_T/\overline{\text{RUL}}_T$, which measures conservatism, that is the cost of
unnecessarily early replacement. With the one-sided score $V=\hat y(X)-Y$ the certified set is
$\{y\ge\hat y(X)-q_T\}$, so a planner acts on the certified fraction of predicted life $1-b$, and
we report it against three fixed bands: $b<0.5$ is
\emph{actionable}, more than half the predicted life being certified; $0.5\le b<1$ is
\emph{degraded}, a positive but shrinking lead time; and $b\ge1$ is \emph{degenerate}, the certified
lower bound being non-positive on average, so that the interval states only that failure has not
yet occurred and carries no replacement schedule. We also report the \emph{operational bound}
$2(a+L\lVert\Delta\boldsymbol\psi\rVert)$ and
the \emph{diagnostic} verdict. Calibration is \emph{unit-level}: every unit contributes one
score, taken at a single randomly chosen monitoring fraction, so the calibration set is
exchangeable across units. Numbers are averaged over five seeds, and coverage gaps are averaged
over (pair, seed) after the clip at zero, so an averaged gap is not the gap of the averaged
coverage; we keep the clipped convention throughout as the conservative reading. Three acceptance
criteria were fixed before running: (1) naive conformal miscovers while SCC recovers at $\eta=0$; (2) the a-priori
bound, with $L$ fit on a subset of (pair, $\eta$) configurations, upper-bounds the gap on
\emph{held-out} configurations; and (3) the coverage gap and the back-off degrade monotonically
with $\eta$. The three bands above are a reporting scale for the envelope, not a fourth criterion.

\subsection{Second case study: an externally authored turbofan fleet}\label{sec:cmapss}
The testbed above is ours, so its departure from similitude is a knob we set. To test the method
where the physics was authored elsewhere we add the NASA C-MAPSS turbofan
datasets~\cite{saxena2008cmapss,saxena2008damage}. FD002 contains 260 run-to-failure engines and
FD004 contains 249, and every engine visits all six operating regimes of the fleet, so each regime
carries roughly 260 units against the six bearings per condition available in the FEMTO benchmark,
a $43\times$ improvement in replication per condition.

The dimensionless coordinate is again derived rather than fitted. Gas-turbine performance collapses
onto referred (corrected) parameters built from the stagnation ratios $\theta=T_t/T_{\mathrm{std}}$
and $\delta_{\mathrm{p}}=P_t/P_{\mathrm{std}}$ (subscripted to keep it distinct from the scalar
departure $\delta$ of Section~\ref{sec:res-cert}), which follow from the recorded altitude, Mach
number and throttle setting through the standard atmosphere and the isentropic stagnation
relations; temperatures are referred by $\theta$, pressures by $\delta_{\mathrm{p}}$, speeds by
$\sqrt{\theta}$ and flows by $\delta_{\mathrm{p}}/\sqrt{\theta}$. The reduction is self-checking:
the sea-level
static regime returns $\theta=1.0000$ and $\delta_{\mathrm{p}}=0.9999$, and the implied damage clock
$\delta_{\mathrm{p}}/\sqrt{\theta}$ spans $3.5\times$ across the six regimes, so the regimes are far
apart in exactly the sense the method cares about.

Degradation is not the operating point but the departure from it, so the coordinate in which SCC
calibrates here is the deviation of each referred sensor from its regime's healthy baseline, the
standard construction of gas-path analysis. That baseline is taken over the first 20 cycles of
operation at the target regime, which a commissioned asset has, and uses no target failure data, so
the method's central claim is preserved. Engines are split disjointly into thirds for predictor
training, calibration and test, because every engine visits every regime and sharing engines would
leak dependence across the split; calibration takes one snapshot per engine per regime, matching the
unit-level scheme of Section~\ref{sec:metrics}. A ridge predictor is fitted on the source regime and
deployed on the target across all 30 ordered regime pairs at $\alpha=0.10$, with a random forest as
a capacity check.

Two departure quantities are computable before any target failure is observed, and both are tested
as candidates for $\boldsymbol\psi$: the distance between regimes in the referred ambient
coordinates $(\log\theta,\log\delta_{\mathrm{p}})$, and the distance between the regimes' healthy
gas-path signatures.
% ---------- end sections/04_experimental_design ----------
% ---------- begin sections/05_results ----------
\section{Results}\label{sec:results}
\subsection{Does SCC restore coverage under regime transfer?}\label{sec:res-cov}
At $\eta=0$, calibrating on one temperature and deploying on another, naive conformal miscovers
badly: coverage collapses to 0.572 when calibrating on short-lived and deploying on long-lived
units (a gap of 0.328), and over-covers wastefully in the reverse direction, reaching 1.000. In
maintenance terms, a 0.328 gap means roughly a third of the intervals promised at 90\% reliability
fail to contain the true life, an unsafe basis for a replacement decision. SCC holds every
calibrate--deploy pair between 0.891 and 0.908, so no pair falls more than 0.009 below the 0.90
target, for a mean coverage gap of 0.011 (Fig.~\ref{fig:cov}), and the non-exchangeable bound
$2\,\dtv(\tilde{\mathcal V})$ contains the SCC gap on every pair. The dimensionless transform removes
the $10.2\times$ scale difference exactly, the behaviour the zero-departure case of~\eqref{eq:cert}
predicts.

\begin{figure}[t]\centering
\includegraphics[width=\linewidth]{fig1_coverage.pdf}
\caption{Coverage under unit-level calibration. (a) Catalyst testbed at $\eta=0$ across
calibrate$\to$deploy temperature pairs: naive either miscovers (0.572) or over-covers (1.000),
and discrete-mode conformal, which has no calibration set at a deployment condition without
failure history and falls back to the nearest mode, inherits the same directional failure, so its
value depends on the deployment condition alone. SCC tracks the 0.90 target (0.891--0.908).
(b) The same comparison on C-MAPSS over all 30 ordered regime pairs of each dataset, group means
marked: naive coverage is bimodal, collapsing to zero on 20 of the 30 pairs and over-covering on
most of the rest, while SCC stays within 0.68--0.97.}\label{fig:cov}
\end{figure}

\subsection{How does the guarantee degrade as physics departs from similitude?}\label{sec:res-deg}
As the unmodelled secondary channel is weighted up, the SCC coverage gap rises monotonically (0.011
at $\eta=0$ to 0.097 at $\eta=2$) and so does the interval back-off, from $b=0.283$ through 0.437 at
$\eta=0.5$ and 0.611 at $\eta=1$ to 1.011 at $\eta=2$ (Fig.~\ref{fig:dep}). Read against the bands
of Section~\ref{sec:metrics} this is an operating envelope rather than a smooth decline: SCC is
actionable out to $\eta=0.68$, where $b$ reaches 0.5, degraded from there to $\eta=1.97$, and
degenerate beyond it. At $\eta=2$ the certified fraction of predicted life is $-0.011$. Coverage is
still certified at that point, the gap being 0.097, but the interval then states only that failure
has not yet occurred and carries no replacement schedule, so $\eta=2$ marks the edge of the useful
envelope rather than a comfortable operating point. We report the envelope because the alternative,
quoting the coverage number alone, leaves a healthy-looking guarantee attached to an interval a
planner cannot act on.

\begin{figure}[t]\centering
\includegraphics[width=\linewidth]{fig2_departure.pdf}
\caption{Under a controlled departure $\eta$ from similitude, the SCC coverage gap (left axis) and
the relative interval back-off (right axis) degrade monotonically. Shading marks the actionability
bands of Section~\ref{sec:metrics}: the actionable band ($b<0.5$) ends at $\eta=0.68$ and
degeneracy ($b\ge1$) begins at $\eta=1.97$, where the certified lower bound is no longer positive
on average.}\label{fig:dep}
\end{figure}

\subsection{Is the coverage bound reliable before deployment?}\label{sec:res-cert}
With the intercept $a$ and slope $L$ fit on 60\% of (pair, $\eta$) configurations, the physical
departure $\delta=\eta\lvert h(S)-h(T)\rvert$, the scalar $m=1$ instance of $\lVert\Delta\boldsymbol\psi\rVert$,
predicts the dimensionless score-distribution mismatch
on the held-out 40\% with $R^2=0.773$ (correlation 0.947), and the resulting a-priori bound
$2(a+L\delta)$ upper-bounds the measured coverage gap on \emph{100\%} of held-out configurations
(Fig.~\ref{fig:cert}). Validity is not sharpness, and the margin says which one this is: across
held-out configurations the mean certified bound is 0.384 against a mean measured gap of 0.046, a
margin of 0.339 (range 0.268--0.552), so the certificate is conservative by roughly a factor of
eight. The looseness sits in the intercept rather than the physical slope, since at $\delta=0$ the
bound is $2a=0.278$ against a measured gap of 0.011. That intercept is the finite-sample floor of
the plug-in total-variation estimator and decays as $n^{-1/2}$ in the number of calibration units
(0.424 at 50 units to 0.096 at 800, a log-log slope of $-0.513$ against the $-0.500$ predicted), so
the certificate tightens as the calibration fleet grows. A second factor of two is irreducible: the
swap argument of Lemma~\ref{lem:barber} bounds the perturbation of the calibration set by twice the
total-variation distance between the source and target score laws, and no amount of calibration
data removes it. For deployment this means the worst-case
coverage shortfall can be quoted from physical parameters before the target asset is run, as a
worst-case planning bound rather than a sharp estimate.

\begin{figure}[t]\centering
\includegraphics[width=\linewidth]{fig3_certificate.pdf}
\caption{Held-out measured coverage gap vs.\ the physical departure $\delta$, with the a-priori
bound $2(a+L\delta)$ fit on disjoint configurations. The bound holds on 100\% of held-out points
and carries a mean margin of 0.339, so the shaded slack is the conservatism discussed in the
text.}\label{fig:cert}
\end{figure}

\subsection{Robustness across operating and tuning choices}\label{sec:res-rob}
The acceptance criteria survive across the departure's temperature-dependence ($E_2=95$--150~kJ/mol
at matched magnitude), its magnitude (15--60\% perturbation), and the coverage target
($\alpha=0.05/0.10/0.20$); Table~\ref{tab:robust}. At every setting naive conformal miscovers
($\approx0.13$), SCC recovers ($\approx0.011$), the a-priori bound holds on 100\% of held-out
configurations while carrying a margin of 0.30--0.50, the same conservatism reported in
Section~\ref{sec:res-cert}, the gap stays monotone, and its magnitude orders with the steepness and
size of the departure. The behaviour is a property of the method rather than of a single tuning.

\begin{table}[t]\centering\footnotesize
\setlength{\tabcolsep}{9pt}\renewcommand{\arraystretch}{1.2}
\caption{Robustness of the acceptance criteria across the departure's temperature-dependence, its
magnitude, and the coverage level, under unit-level calibration. $g_0$: coverage gap at $\eta=0$;
$g_{\max}$: SCC coverage gap at $\eta=2$; Margin: mean of (certified bound $-$ measured gap) over
held-out configurations, so a positive value means valid and a large value means loose; Bound\,\%:
fraction of held-out configurations on which the a-priori bound holds.}\label{tab:robust}
\begin{tabular}{@{}llccccc@{}}\toprule
Perturbation axis & Value & Naive $g_0$ & SCC $g_0$ & SCC $g_{\max}$ & Margin & Bound\,\% \\\midrule
$E_2$ (kJ/mol) & 95   & 0.132 & 0.011 & 0.043 & 0.302 & 100 \\
               & 150  & 0.132 & 0.011 & 0.218 & 0.502 & 100 \\
Perturbation   & 15\,\% & 0.132 & 0.011 & 0.055 & 0.302 & 100 \\
               & 60\,\% & 0.132 & 0.011 & 0.131 & 0.400 & 100 \\
Coverage $\alpha$ & 0.05 & 0.138 & 0.008 & 0.071 & 0.372 & 100 \\
               & 0.20 & 0.103 & 0.011 & 0.127 & 0.323 & 100 \\\bottomrule
\end{tabular}
\end{table}

The result also survives a change of base predictor, which is the sharper form of the question:
a more flexible learner might be expected to degrade the guarantee. The predictor is developed on
the source condition and deployed on the target, the naive baseline fitting it in raw coordinates
and SCC fitting the same model class in the dimensionless ones
(Table~\ref{tab:basemodel}). The risk is real and it falls entirely on the unscaled baseline: the
naive gap at $\eta=0$ worsens from 0.127 for the closed-form extrapolator to 0.449 for a random
forest, a factor of $3.5$, because a learner fitted flexibly to one operating regime fails harder
when deployed on another. SCC is flat at 0.007 to 0.014 across all four model classes, and under
departure the more flexible learners are the better ones (0.046 for the neural network against
0.102 for the closed form at $\eta=2$), since a model fitted in dimensionless coordinates absorbs
part of the residual departure. Conservatism does not degrade with capacity.

Two qualifications belong with that claim. The drop-in property holds provided the learned
predictor is fitted in the dimensionless coordinates; a model fitted in raw coordinates cannot be
repaired by calibration alone, because no conformal layer can fix a point predictor that does not
transfer. And the ordering of the naive column is a property of the testbed rather than a law: on
C-MAPSS it runs the other way (Section~\ref{sec:res-cmapss}), where the raw sensors encode the
operating regime itself, so a flexible learner can partly infer the regime and compensate, while on
the testbed the raw signal carries no such marker. What is invariant across both datasets is that
SCC is the better arm in every row.

The reviewer's question extends to the certificate, so we fit it separately for each predictor
class over the same configurations. Validity is unaffected: the a-priori bound holds on 100\% of
held-out configurations for all four, with margins in the narrow band 0.30 to 0.36, so its
conservatism does not grow with capacity. Its \emph{mechanism} does weaken. The departure model
explains the score-distribution mismatch well only for the closed-form predictor (held-out
$R^2=0.90$ against $-2.8$ to $-0.02$ for the three learned ones), so with a learned predictor the
bound holds on conservatism rather than because $\delta$ tracks the mismatch. The residual there is
dominated by model error rather than by physical departure, which is the same pattern the C-MAPSS
certificate shows in Section~\ref{sec:res-cmapss} and a further reason the certificate is claimed
narrowly.

\begin{table}[t]\centering\footnotesize
\setlength{\tabcolsep}{10pt}\renewcommand{\arraystretch}{1.2}
\caption{Coverage gap by base predictor, trained on the source condition and deployed on the
target, under unit-level calibration. Naive fits the predictor in raw coordinates; SCC fits the same
model class in the dimensionless ones. The last two columns fit the a-priori certificate separately
for each predictor over the same configurations: the margin is the mean of (bound $-$ measured gap)
on held-out configurations and $R^2$ is the held-out fit of the departure model. The bound holds on
100\% of held-out configurations for every predictor class.}\label{tab:basemodel}
\begin{tabular}{@{}lccccrr@{}}\toprule
& \multicolumn{2}{c}{$\eta=0$} & \multicolumn{2}{c}{$\eta=2$}
& \multicolumn{2}{c}{Certificate} \\\cmidrule(lr){2-3}\cmidrule(lr){4-5}\cmidrule(lr){6-7}
Base predictor & Naive & SCC & Naive & SCC & Margin & $R^2$ \\\midrule
Closed-form physics & 0.127 & 0.014 & 0.335 & 0.102 & 0.361 & 0.90 \\
Ridge regression    & 0.275 & 0.008 & 0.317 & 0.061 & 0.323 & $-2.76$ \\
Random forest       & 0.449 & 0.014 & 0.450 & 0.085 & 0.299 & $-0.02$ \\
Neural network      & 0.346 & 0.007 & 0.483 & 0.046 & 0.325 & $-0.03$ \\\bottomrule
\end{tabular}
\end{table}

\subsection{Does the method transfer to a fleet we did not author?}\label{sec:res-cmapss}
On C-MAPSS the calibration transfers and the a-priori certificate does not. The two halves scope the
contribution differently, so we report them separately.

Regime-blind conformal fails harder on the real fleet than on the testbed: calibrating on one regime
and deploying on another loses a mean 0.590 of coverage across the 30 ordered pairs of FD002.
Referring the sensors to standard day is not by itself enough and still loses 0.550; calibrating in
the gas-path deviation coordinate recovers coverage to within 0.060 of the 0.90 target, a
$9.8\times$ reduction in the gap, and FD004 replicates it at 0.567 to 0.058
(Table~\ref{tab:cmapss}). The mechanism therefore survives contact with an externally authored,
well-powered dataset, still without target failure data. The size of the naive failure does depend
on the predictor: a random forest on raw sensors can partly infer the regime from the sensor levels
themselves and loses only 0.128 on FD002, but SCC improves on it there too (0.053), so the recovery
is not an artefact of a weak baseline.

The certificate is a different matter. Neither departure candidate identifies the residual gap: the
ambient referred distance correlates 0.289 with it on FD002 and 0.179 on FD004, and the healthy
gas-path signature distance correlates $-0.122$ and $-0.086$. Held out, both regressions explain
nothing, with $R^2$ between $-0.074$ and 0.022. The gap is dominated instead by an offset that the
departure term does not touch, the fitted intercepts being 0.026 to 0.058 against mean gaps of 0.060
and 0.058. The bound consequently holds on 67 to 83\% of held-out regime pairs, against 100\% on the
testbed, where the departure predicts the score-distribution mismatch the bound is built from with
$R^2=0.773$.

We report this as the boundary of the contribution rather than as a failure of it. The dimensionless
calibration, which is the part a practitioner deploys, is now validated on real, externally authored
data with about 260 units per regime and replicated across two datasets. The a-priori certificate
requires $\boldsymbol\psi$ to be identified from the governing degradation law, which the catalyst
testbed supplies and a turbofan gas path does not, at least not through either candidate tested
here.~\ref{app:bearing} works that identification through for rolling-contact fatigue, where
ISO~281 enumerates the residual groups directly, and the contrast with the gas path is the clearest
statement we can give of where the certificate applies. Where $\boldsymbol\psi$ is unavailable the
runtime diagnostic of Section~\ref{sec:diagnostic} is the deployment safeguard, and
Section~\ref{sec:res-diag} shows what it does when the evidence is too thin to support the
premise.

\begin{table}[t]\centering\footnotesize
\setlength{\tabcolsep}{7pt}\renewcommand{\arraystretch}{1.2}
\caption{C-MAPSS, 30 ordered regime pairs per dataset, target coverage 0.90, unit-level calibration
on engine-disjoint splits. Top: coverage gap for the regime-blind baseline and for SCC calibrated in
the gas-path deviation coordinate. Bottom: the a-priori certificate with $\boldsymbol\psi$ taken in
turn as each candidate; $R^2$ is on held-out regime pairs, and the same certificate holds on 100\%
of held-out configurations on the catalyst testbed.}\label{tab:cmapss}
\begin{tabular}{@{}llccc@{}}\toprule
Dataset & Base model & Naive & SCC & Improvement \\\midrule
FD002 (260 engines) & ridge & 0.590 & 0.060 & $9.8\times$ \\
FD004 (249 engines) & ridge & 0.567 & 0.058 & $9.8\times$ \\
FD002 (260 engines) & random forest & 0.128 & 0.053 & $2.4\times$ \\
FD004 (249 engines) & random forest & 0.123 & 0.033 & $3.7\times$ \\\midrule
Dataset & $\boldsymbol\psi$ candidate & corr & held-out $R^2$ & Bound\,\% \\\midrule
FD002 & ambient referred distance & 0.289 & 0.022 & 67 \\
FD002 & healthy gas-path signature & $-0.122$ & $-0.037$ & 67 \\
FD004 & ambient referred distance & 0.179 & $-0.010$ & 75 \\
FD004 & healthy gas-path signature & $-0.086$ & $-0.074$ & 83 \\\bottomrule
\end{tabular}
\end{table}

\subsection{Comparison with discrete-mode conformal}\label{sec:res-mondrian}
The natural alternative to a physical scale is to treat each operating condition as a discrete mode
and calibrate within it~\cite{javanmardi2023conformal}. We give that baseline the most favourable
reading available, supplying the three testbed temperatures as known modes, and separate the two
cases that differ operationally (Table~\ref{tab:mondrian}).

When the target condition has its own run-to-failure history, discrete-mode calibration is
near-oracle and, once the departure is large, better than SCC (0.009 against 0.045 at $\eta=1$). SCC
matches it under exact similitude (0.009 against 0.011) while using no target failure data at all.
We state the implication plainly: where failure data already exists at the deployment condition, the
discrete remedy is the right tool and the physical scale is not needed.

The second case is the one SCC exists for. When the target condition has no failure history the
discrete method has no calibration set for that mode, must fall back to the nearest mode, and lands
close to regime-blind pooling (0.078 to 0.138 against 0.093 to 0.159 pooled). SCC holds at 0.009 to
0.040, a gap smaller by $3.5\times$ to $8.7\times$, because the physical scale transports the
calibration instead of requiring it to be re-earned at every condition. A newly commissioned unit, a
new operating point, or any condition the fleet has not yet run to failure is always in this case.

\begin{table}[t]\centering\footnotesize
\setlength{\tabcolsep}{10pt}\renewcommand{\arraystretch}{1.2}
\caption{Coverage gap against discrete-mode conformal, unit-level calibration, mean over target
conditions and five seeds. \emph{Target history} indicates whether the deployment condition has its
own run-to-failure data; pooled is regime-blind calibration over all conditions and does not apply
when the target has its own data.}\label{tab:mondrian}
\begin{tabular}{@{}llccc@{}}\toprule
$\eta$ & Target history & Discrete-mode & SCC & Pooled \\\midrule
0.0 & available   & 0.011 & 0.009 & n/a \\
0.5 & available   & 0.017 & 0.031 & n/a \\
1.0 & available   & 0.009 & 0.045 & n/a \\\midrule
0.0 & unavailable & 0.078 & \textbf{0.009} & 0.093 \\
0.5 & unavailable & 0.109 & \textbf{0.025} & 0.128 \\
1.0 & unavailable & 0.138 & \textbf{0.040} & 0.159 \\\bottomrule
\end{tabular}
\end{table}

\subsection{What happens when the assumed physics is wrong?}\label{sec:res-misspec}
The scale is built from an assumed activation energy, so the method has to be tested against that
assumption being wrong. We perturb the assumed $E_1$ while holding $k(T_{\mathrm{ref}})$ fixed, so
only the between-condition behaviour of $\sigma$ is distorted, which is the part similitude depends
on (Table~\ref{tab:misspec}).

Two cases separate sharply. When only the scale is wrong and the predictor is unaffected, the
tolerance is wide: the SCC gap stays at or below 0.05 out to $\lvert\Delta E/E\rvert=31\%$, remains
better than not scaling at all out to $+62\%$, and becomes worse than the unscaled baseline only at
$+100\%$ (0.139 against 0.132). When a single wrong model drives both the predictor and the scale,
which is the realistic case, the method is much more fragile: SCC is already worse than the unscaled
baseline at $-10\%$ (0.072 against 0.061) and clearly worse at $-19\%$ (0.157 against 0.030).

The mitigation is that the dangerous regime is detectable before deployment, from source data alone.
The dimensionless-life invariance check returns \emph{holds} only at the correct scale ($g=0.967$,
interval $[0.92,1.01]$) and \emph{violated} at every misspecification tested, including the $-10\%$
point at which the coupled case first turns against the method ($g=1.215$, $[1.16,1.27]$) and the
$+10\%$ point ($g=0.769$, $[0.73,0.81]$). This is a capability the submitted manuscript did not
claim: the runtime check tests the invariance of the dimensionless life, so it is sensitive to
misspecification of the scale itself as well as to a departure from similitude. The failure mode the
reviewer identifies is real, and it is not silent.

\begin{table}[t]\centering\footnotesize
\setlength{\tabcolsep}{7pt}\renewcommand{\arraystretch}{1.2}
\caption{Misspecified activation energy. \emph{Scale only}: the assumed $E_1$ enters the
dimensionless scale but not the predictor, so the naive baseline is 0.132 at every row by
construction. \emph{Coupled}: the same wrong model drives predictor and scale. $g$ and its bootstrap
interval are the runtime invariance check computed from source data only, with \emph{holds} returned
when 1 lies inside the interval.}\label{tab:misspec}
\begin{tabular}{@{}lccccc@{}}\toprule
& Scale only & \multicolumn{2}{c}{Coupled} & \multicolumn{2}{c}{Invariance check} \\
\cmidrule(lr){2-2}\cmidrule(lr){3-4}\cmidrule(lr){5-6}
$\Delta E/E$ & SCC & Naive & SCC & $g$ & 95\% interval \\\midrule
$-50\%$  & 0.076 & 0.326 & 0.367 & 3.039 & [2.90, 3.19] \\
$-31\%$  & 0.047 & 0.246 & 0.275 & 1.978 & [1.89, 2.07] \\
$-19\%$  & 0.029 & 0.030 & 0.157 & 1.485 & [1.42, 1.56] \\
$-10\%$  & 0.016 & 0.061 & 0.072 & 1.215 & [1.16, 1.27] \\
$0\%$    & 0.011 & 0.132 & 0.011 & 0.967 & [0.92, 1.01] \\
$+10\%$  & 0.020 & 0.211 & 0.055 & 0.769 & [0.73, 0.81] \\
$+19\%$  & 0.036 & 0.269 & 0.142 & 0.629 & [0.60, 0.66] \\
$+62\%$  & 0.105 & 0.381 & 0.401 & 0.231 & [0.22, 0.24] \\
$+100\%$ & 0.139 & 0.393 & 0.443 & 0.098 & [0.09, 0.10] \\\bottomrule
\end{tabular}
\end{table}

\subsection{Knowing when not to trust the method: the design envelope and a thin experimental
benchmark}\label{sec:res-diag}
A diagnostic that abstains is useful only if the data environment it needs is attainable, so we map
that first. Under exact similitude, where \emph{holds} is the correct answer, the fraction of trials
returning it depends on the units per condition and on the unit-to-unit life scatter
(Table~\ref{tab:power}). The requirement is modest: three units suffice at a scatter of 0.10, six at
0.25 and ten at 0.50, which is within reach of an ordinary industrial fleet, and the C-MAPSS fleet
at roughly 260 units per regime sits far inside the feasible region. Only at a scatter of 0.80 does
the requirement rise to about forty units. The map should be read as a floor on fleet size rather
than as a power curve approaching one: the reported verdict is the worst of three pairwise interval
tests at a nominal 95\% each, so even under exact similitude the achievable rate is about
$0.95^3\approx0.86$, which is where the larger fleets settle. A multiplicity-corrected verdict is
the obvious refinement and would raise that ceiling.

\begin{table}[t]\centering\footnotesize
\setlength{\tabcolsep}{12pt}\renewcommand{\arraystretch}{1.2}
\caption{Fraction of trials in which the diagnostic returns the correct \emph{holds} verdict under
exact similitude, by units per condition and unit-to-unit life scatter $\sigma_{\ln A}$. Forty
trials per cell, so each entry carries a standard error of about 0.06.}\label{tab:power}
\begin{tabular}{@{}lcccc@{}}\toprule
& \multicolumn{4}{c}{Life scatter $\sigma_{\ln A}$} \\\cmidrule(lr){2-5}
Units per condition & 0.10 & 0.25 & 0.50 & 0.80 \\\midrule
3   & 0.90 & 0.75 & 0.20 & 0.05 \\
6   & 0.97 & 0.90 & 0.50 & 0.00 \\
10  & 0.95 & 0.88 & 0.82 & 0.17 \\
20  & 0.97 & 0.95 & 0.85 & 0.68 \\
40  & 1.00 & 0.95 & 0.85 & 0.93 \\
80  & 0.97 & 0.82 & 0.90 & 0.88 \\
160 & 0.90 & 0.78 & 0.88 & 0.82 \\\bottomrule
\end{tabular}
\end{table}

On synthetic data the runtime diagnostic discriminates correctly: it returns \emph{holds} when lives
obey L10, \emph{violated} when they do not (with adequate units), and \emph{indeterminate} when the
within-condition spread swamps the structural signal (Fig.~\ref{fig:diag}). The FEMTO bearing
benchmark~\cite{nectoux2012pronostia} is an accelerated laboratory test platform, not field data
from an operating fleet, and it falls in the last category: with roughly six bearings per
condition and a $7\times$ within-condition life spread, the bootstrap CI on $g$ spans a factor of
2.6--4.4, far too wide to resolve the $1.63\times$ life ratio that L10 predicts between its load/speed
conditions (derived for these conditions in~\ref{app:bearing}). The diagnostic therefore returns
\emph{indeterminate} and SCC declines its a-priori
certificate, deferring to data-driven conformal. This is an honest applicability limit rather than a
failure, and it is consistent with the dataset authors' report that L10 does not govern FEMTO lives;
it also shows the diagnostic doing its safety job, refusing to certify on evidence too weak to
support the premise. The map also predicts this outcome rather than excusing it: at six units per
condition with a spread comparable to the two rightmost columns of Table~\ref{tab:power},
\emph{holds} is returned in 0.50 to 0.00 of trials and \emph{indeterminate} in 0.45 to 0.90, so
the verdict follows from the data environment rather than from anything peculiar to the dataset.

\begin{figure}[t]\centering
\includegraphics[width=\linewidth]{fig4_diagnostic.pdf}
\caption{Diagnostic verdict via bootstrap CIs on $g$ (median-life ratio normalised by the L10
prediction; $g{=}1$ under similitude). Synthetic similitude $\to$ holds; powered violation $\to$
violated; FEMTO-like (few units, large spread) $\to$ indeterminate.}\label{fig:diag}
\end{figure}
% ---------- end sections/05_results ----------
% ---------- begin sections/06_discussion ----------
\section{Discussion}\label{sec:discussion}
\subsection{What the guarantee provides in practice}\label{sec:disc-buys}
SCC quantifies the reliability of a transferred RUL interval from physical parameters before any
target failure data exists. When the runtime diagnostic confirms similitude (\emph{holds}),
arbitrarily large differences in the dominant operating parameters, the $10.2\times$ life range in
the testbed, cost nothing in coverage, because the dimensionless scale absorbs them. When similitude
is incomplete, the certificate degrades by the physical departure $\lVert\Delta\boldsymbol\psi\rVert$
and the diagnostic flags the regime. This is a different bargain from weighted
conformal~\cite{tibshirani2019conformal}: SCC trades the need for target failure data, which a new
asset lacks, for a falsifiable physical assumption that the diagnostic makes checkable.

\subsection{Deployment and fault-tolerant use}\label{sec:disc-deploy}
The two by-products of the method are what make it usable on a real asset. The a-priori bound lets a
planner quote a worst-case coverage shortfall, and therefore a worst-case reliability of the
maintenance decision, before the target unit is commissioned. The runtime diagnostic acts as an
online monitor: when its premise is unsupported it returns \emph{indeterminate} or \emph{violated},
and the system falls back to conservative data-driven conformal rather than issuing an interval whose
confidence is silently wrong. The result is announced degradation rather than an unannounced loss
of coverage: where the method weakens, the weakening is visible in the reported efficiency and in
the diagnostic verdict, which is the property a safety-relevant prognostic function needs.

\subsection{Limitations and validation roadmap}\label{sec:disc-limits}
Four limitations bound the claims. (i) The guarantee is conditional on the physical hypotheses
H1--H2 (\ref{app:proof}), which are falsifiable rather than proven. (ii) The certificate bounds the
gap for the modelled groups; unmodelled physics contributes a residual that the departure study
quantifies and the diagnostic detects, but that must be bounded a priori to use the certificate under
incomplete similitude. (iii) The evidence is graded rather than uniform: the coverage recovery is
demonstrated on a controlled simulation and replicated on an externally authored turbofan fleet,
but the a-priori certificate is validated only where the governing law identifies
$\boldsymbol\psi$, and the one experimental benchmark, an accelerated laboratory test platform
rather than field data, returns \emph{indeterminate} by design rather than a positive result.
(iv) The bound is conservative.
The decisive next step is validation on adequately powered run-to-failure data from an operating
fleet, where the diagnostic returns \emph{holds} and the certificate can be checked against
measured cross-regime coverage; such data is the goal of ongoing work on an operating industrial
process.
% ---------- end sections/06_discussion ----------
% ---------- begin sections/07_conclusions ----------
\section{Conclusions}\label{sec:conclusions}
We presented Similarity-Calibrated Conformal prediction, a deployable method that delivers reliable
remaining-useful-life intervals when an asset is operated in a regime different from the one its model
was calibrated on. By calibrating the conformal quantile in a physics-derived dimensionless
coordinate, SCC restores coverage under operating-regime transfer with no target failure data,
filling the gap between dimensionless transfer learning, which lacks a coverage guarantee, and
weighted conformal, which needs target data.

The evidence supports two claims of different strength, and we separate them. The dimensionless
calibration is the deployable part and is validated twice. On a controlled catalyst-deactivation
reliability testbed it recovers coverage where regime-blind conformal fails, degrades monotonically
and within a stated actionability envelope under a controlled departure from similitude, and is
robust across the structure and magnitude of the departure, the coverage level and the choice of
base predictor. On 509 externally authored turbofan engines it cuts the cross-regime coverage gap
from 0.590 to 0.060, replicated across two datasets. The a-priori bound is the narrower claim. It
holds on every held-out configuration of the testbed but with a mean margin of 0.339, so it is a
planning bound rather than a sharp estimate, and it requires the governing degradation law to
identify the similitude-departure quantity. ISO~281 supplies that identification for rolling-contact
fatigue (see~\ref{app:bearing}); no comparable enumeration exists for a turbofan gas path, and
neither candidate tested there predicted the residual gap. We therefore claim the certificate for
assets whose degradation law supplies $\boldsymbol\psi$, and the calibration and the runtime
diagnostic more broadly.

On the FEMTO experimental benchmark, an accelerated laboratory test platform rather than field data,
the diagnostic correctly places that dataset outside the method's envelope, demonstrating the
intended fault-tolerant behaviour; the design map of Section~\ref{sec:res-diag} shows that verdict
was predictable from six bearings per condition and a sevenfold life spread rather than peculiar to
the dataset. The decisive test remains validation on adequately powered run-to-failure field data
from an operating process, which is the focus of continuing work.
% ---------- end sections/07_conclusions ----------

\section*{Data and code availability}
The SCC implementation, the deactivation testbed, the diagnostic, and scripts reproducing every
figure and table are available at \texttt{github.com/beebzy-droid/IPIS}.

\section*{CRediT authorship contribution statement}
\textbf{Bien Busico:} Conceptualization, Methodology, Software, Formal analysis, Validation,
Writing -- original draft.

\section*{Declaration of generative AI use}
During preparation the author used a generative AI assistant for derivation checking, code
scaffolding, and editing. The author reviewed and verified all content and takes full
responsibility for it.


\appendix
% appendix floats restart per appendix section (A.1, B.1, ...)
\setcounter{table}{0}\setcounter{figure}{0}
\renewcommand{\thetable}{\Alph{section}.\arabic{table}}
\renewcommand{\thefigure}{\Alph{section}.\arabic{figure}}
% ---------- begin sections/08_appendix ----------
\section{Formal coverage guarantee and proof}\label{app:proof}
This appendix gives the formal statement underlying Eq.~\eqref{eq:cert}: the dimensionless reduction,
the regularity hypotheses, the non-exchangeable bound used as a lemma, and the SCC certificate with
its proof and two-regime corollary.

\begin{definition}[Dimensionless reduction]\label{def:pi}
The Buckingham $\Pi$ theorem yields dimensionless groups, partitioned into (i) \emph{scale-setting}
groups, used to form a characteristic score scale $\sigma(\theta_c)>0$ such that the dimensionless
score $\tilde V=V/\sigma(\theta_c)$ has a condition-invariant law under exact similitude, and (ii) a
\emph{similitude-departure parameter} $\boldsymbol\psi(\theta_c)\in\mathbb R^m$ collecting residual
groups the scale does not absorb. $\boldsymbol\psi$ is constant across conditions when the
dimensionless model is complete. Let $\lawt c$ denote the law of $\tilde V$ under $\mathbb P_c$.
\end{definition}

\begin{assumption}[Lipschitz similitude departure]\label{ass:lip}
There exists $L<\infty$ with
$\dtv(\lawt c,\lawt{c'})\le L\lVert\boldsymbol\psi(\theta_c)-\boldsymbol\psi(\theta_{c'})\rVert$ for
all $c,c'$.
\end{assumption}
The exact-similitude boundary case
($\boldsymbol\psi\equiv\mathrm{const}\Rightarrow\lawt c=\lawt{c'}$, hypothesis H1) holds
\emph{irrespective of how far apart the scale-setting parameters are}; the global Lipschitz form
(hypothesis H2) is the regularity needed away from exact similitude. H1--H2 are physical hypotheses,
not theorems, and the diagnostic of Section~\ref{sec:diagnostic} tests them.

\begin{lemma}[\cite{barber2023conformal}]\label{lem:barber}
For split conformal with weights $w_i$ and pre-fitted scores, the coverage gap is at most
$(1+\sum_i w_i)^{-1}\sum_i w_i\,\dtv(R(Z),R(Z^i))$, where $R(Z)$ is the score vector and $R(Z^i)$
swaps the test point with calibration point $i$. For independent data,
$\dtv(R(Z),R(Z^i))\le 2\,\dtv(Z_i,Z_{n+1})$.
\end{lemma}

\begin{theorem}[SCC certificate]\label{thm:scc}
Under Assumption~\ref{ass:lip}, with calibration i.i.d.\ from $S$ and target
$(X_{n+1},Y_{n+1})\sim\mathbb P_T$ independent of calibration,
\[
\mathbb P\big(Y_{n+1}\in\mathcal C(X_{n+1})\big)\ge
1-\alpha-\tfrac{2n}{n+1}L\lVert\boldsymbol\psi(\theta_S)-\boldsymbol\psi(\theta_T)\rVert
\ge 1-\alpha-2L\lVert\boldsymbol\psi(\theta_S)-\boldsymbol\psi(\theta_T)\rVert.
\]
\end{theorem}
\begin{proof}
Coverage is the event $\{\tilde V_{n+1}\le\tilde q\}$ with $\tilde V_{n+1}\sim\lawt T$. Were the
dimensionless scores exchangeable (all $\lawt S$), split conformal would give coverage $\ge1-\alpha$.
Apply Lemma~\ref{lem:barber} with $w_i\equiv1$ (denominator $n+1$). The pre-fitted scores are
independent scalars; swapping calibration point $i$ with the test point exchanges
$\tilde V_i\sim\lawt S$ and $\tilde V_{n+1}\sim\lawt T$ while the other coordinates are unchanged. The
two joint laws are product measures differing in two independent coordinates, so by subadditivity of
total variation $\dtv(R(Z),R(Z^i))\le\dtv(\lawt S,\lawt T)+\dtv(\lawt T,\lawt S)=2\dtv(\lawt S,\lawt T)$.
Summing $n$ terms gives gap $\le\tfrac{2n}{n+1}\dtv(\lawt S,\lawt T)$, and Assumption~\ref{ass:lip}
bounds $\dtv(\lawt S,\lawt T)$ by $L\lVert\Delta\boldsymbol\psi\rVert$.
\end{proof}

\begin{corollary}[A-priori certificate, two regimes]\label{cor:apriori}
$\lVert\Delta\boldsymbol\psi\rVert$ is a physical quantity, not a function of target failure data.
\emph{(a)} Under complete similitude ($\boldsymbol\psi$ constant) the bound is $0$ and coverage is
$\ge1-\alpha$ for arbitrary differences in the scale-setting parameters, requiring only the target
scale $\sigma(\theta_T)$. \emph{(b)} Under incomplete similitude, coverage is certified down to
$1-\alpha-2L\lVert\Delta\boldsymbol\psi\rVert$ given a bound on the departure. Neither regime estimates
a density ratio, the property that distinguishes SCC from weighted and covariate-shift conformal.
\end{corollary}

\section{Worked instance: $\sigma$ and $\boldsymbol\psi$ for rolling-contact fatigue}\label{app:bearing}
Definition~\ref{def:pi} splits the dimensionless groups into a scale $\sigma$ and a residual
departure $\boldsymbol\psi$, and the certificate is only as concrete as those two objects. This
appendix carries the split through for the best-standardised degradation law in the reliability
literature, rolling-contact fatigue, and evaluates both on the FEMTO operating
conditions~\cite{nectoux2012pronostia}.

\subsection*{The scale}
The governing quantities are life $t$~[T], dynamic equivalent load $P$~[F], basic dynamic load
rating $C$~[F] and rotational speed $n$~[1/T]. Four quantities in two dimensions give two
independent groups,
\[
\Pi_1=t\,n \quad\text{(revolutions)},\qquad \Pi_2=P/C,
\]
and Lundberg--Palmgren~\cite{lundberg1947dynamic,iso281} is precisely a relation between them,
$\Pi_1=10^{6}\,\Pi_2^{-p}$ with $p=3$ for ball bearings. Solving for life in hours gives the scale
\begin{equation}\label{eq:sigma-bearing}
\sigma=\frac{10^{6}}{60\,n}\left(\frac{C}{P}\right)^{p}\quad[\text{h}],
\end{equation}
so $\sigma$ absorbs load and speed exactly. Table~\ref{tab:bearing-sigma} evaluates
Eq.~\eqref{eq:sigma-bearing} on the three published FEMTO conditions for the 6804 deep-groove ball
bearing of that testbed, using its catalogue rating $C=4$~kN and mean diameter $d_m=26$~mm. The
characteristic life varies by $1.63\times$ across the conditions, which recovers from first
principles the L10 ratio quoted in Section~\ref{sec:res-diag}.

\begin{table}[h]\centering\footnotesize
\setlength{\tabcolsep}{10pt}\renewcommand{\arraystretch}{1.2}
\caption{The scale on the FEMTO conditions, from Eq.~\eqref{eq:sigma-bearing} with $C=4$~kN,
$p=3$. L10 is normalised to condition~1.}\label{tab:bearing-sigma}
\begin{tabular}{@{}lcccc@{}}\toprule
Condition & $n$ [rpm] & $P$ [N] & L10 [normalised] & $\sigma$ [h] \\\midrule
1 & 1800 & 4000 & 1.000 & 9.26 \\
2 & 1650 & 4200 & 0.864 & 8.73 \\
3 & 1500 & 5000 & 0.512 & 5.69 \\\bottomrule
\end{tabular}
\end{table}

\subsection*{The departure parameter}
ISO~281 does not stop at the basic rating life. The modified life is
$L_{nm}=a_1\,a_{\mathrm{ISO}}\,L_{10}$, where the life modification factor
$a_{\mathrm{ISO}}=f(e_C C_u/P,\ \kappa)$ depends on the contamination group $e_C C_u/P$ and the
viscosity ratio $\kappa=\nu/\nu_1$. The arguments of $a_{\mathrm{ISO}}$ are exactly the groups the
basic rating life does not absorb, so for this law the standard itself enumerates
$\boldsymbol\psi$: the lubrication regime, principally $\kappa$, together with the contamination
group. Carrying both makes $\boldsymbol\psi$ two-dimensional, in which case the norm in
Eq.~\eqref{eq:cert} must be fixed alongside $\boldsymbol\psi$ as Section~\ref{sec:guarantee} notes. We
take the scalar $\boldsymbol\psi=\ln\kappa$, the viscosity ratio being multiplicative, because the FEMTO
record reports nothing about contamination. The reference
viscosity is fixed by geometry and speed, $\nu_1=4500\,n^{-0.5}d_m^{-0.5}$ for $n\ge1000$~rpm, while
$\nu$ follows from the lubricant and its \emph{operating temperature} through the Walther relation
of ASTM~D341. Two conditions are in similitude when they share $\kappa$, whatever their loads and
speeds.

Table~\ref{tab:bearing-psi} evaluates $\boldsymbol\psi$ on the same three conditions for an
illustrative ISO~VG~100 grease base oil. The lubricant and its temperatures are illustrative, since
the FEMTO record does not report them; the point is the comparison between the two rows, not the
absolute values.

\begin{table}[h]\centering\footnotesize
\setlength{\tabcolsep}{9pt}\renewcommand{\arraystretch}{1.2}
\caption{The departure parameter on the FEMTO conditions for an illustrative ISO~VG~100 base oil,
with $\nu_1=4500\,n^{-0.5}d_m^{-0.5}$ and $\boldsymbol\psi=\ln\kappa$. The last column is the
largest pairwise departure $\lVert\Delta\boldsymbol\psi\rVert$ entering the
certificate.}\label{tab:bearing-psi}
\begin{tabular}{@{}lcccc@{}}\toprule
Thermal state & $\kappa_1$ & $\kappa_2$ & $\kappa_3$ & $\lVert\Delta\boldsymbol\psi\rVert$ \\\midrule
All conditions at 60\,$^\circ$C            & 1.933 & 1.851 & 1.765 & 0.091 \\
Spread 50 / 60 / 75\,$^\circ$C             & 2.957 & 1.851 & 1.025 & 1.060 \\\bottomrule
\end{tabular}
\end{table}

\subsection*{What this tells a practitioner}
Load and speed differ enough across these conditions to move the characteristic life by
$1.63\times$, and they cost nothing in similitude because $\sigma$ absorbs them exactly. At a matched
thermal state the residual departure is 0.091; a 25\,$^\circ$C spread in operating temperature
raises it to 1.060, roughly $12\times$ larger. For rolling-contact fatigue it is therefore the
thermal and lubrication state, not the load, that breaks similitude and has to be matched or
reported as $\boldsymbol\psi$. That is an actionable instruction for commissioning a fleet, and it
is not the instruction an engineer would guess from the rating-life formula alone.

It also delimits the certificate. Where the governing standard enumerates the residual groups, as
ISO~281 does through $a_{\mathrm{ISO}}$, $\boldsymbol\psi$ is identifiable and the a-priori bound is
applicable. No equivalent enumeration exists for a turbofan gas path, which is the most plausible
reading of the negative result in Section~\ref{sec:res-cmapss}: identifiability of
$\boldsymbol\psi$, rather than the bound itself, is the practical boundary of the method.
% ---------- end sections/08_appendix ----------

\bibliographystyle{elsarticle-num-names}
\bibliography{scc_refs}
\end{document}
